\documentclass[../../main.tex]{subfiles}
\graphicspath{{\subfix{../../image/}}}

\begin{document}

\section*{习题四}\label{section:ch04-习题}

\begin{exercise}\label{exercise:ch04-01}
(1) 证明关于示性函数成立如下公式：
\[
\mathbf{1}_{\bar{A}}=1-\mathbf{1}_A;\qquad
\mathbf{1}_{AB}=\mathbf{1}_A\cdot\mathbf{1}_B;\qquad
\mathbf{1}_{A\cup B}=1-\mathbf{1}_{\overline{AB}}.
\]
(2) 利用示性函数导出概率加法公式~(1.5.6)。

(3) 证明：$\mathbf{1}_{A\cup B\cup C}=\mathbf{1}_A+\mathbf{1}_B+\mathbf{1}_C-\mathbf{1}_{AB}-\mathbf{1}_{BC}-\mathbf{1}_{AC}+\mathbf{1}_{ABC}$。

(4) 证明：$\mathbf{1}_{ABC}=\mathbf{1}_A+\mathbf{1}_B+\mathbf{1}_C-\mathbf{1}_{A\cup B}-\mathbf{1}_{B\cup C}-\mathbf{1}_{A\cup C}+\mathbf{1}_{A\cup B\cup C}$。

(5) 有五个队参加的比赛中，每个队与别的队都比赛一场，若每场比赛参赛双方各有 $50\%$ 赢的机会，试求整个比赛既没有不败的队也没有不胜的队的概率。

(6) 证明 $A$ 与 $B$ 独立的充要条件为 $\mathbf{1}_A$ 与 $\mathbf{1}_B$ 独立。
\end{exercise}

\begin{exercise}\label{exercise:ch04-02}
随机变量 $\mu$ 取非负整数值 $n\geqslant 0$ 的概率为 $p_n=A\dfrac{B^n}{n!}$，已知 $E\mu=a$，试决定 $A$ 与 $B$。
\end{exercise}

\begin{exercise}\label{exercise:ch04-03}
设随机变量 $\xi$ 只取非负整数值，其概率为 $P\{\xi=k\}=\dfrac{a^k}{(1+a)^{k+1}},\ a>0$ 是常数，试求 $E\xi$ 及 $D\xi$。
\end{exercise}

\begin{exercise}\label{exercise:ch04-04}
若事件 $A$ 在第 $i$ 次试验中出现的概率为 $p_i$，设 $\mu$ 是事件 $A$ 在起初 $n$ 次独立试验中的出现次数，试求 $E\mu$ 及 $D\mu$。
\end{exercise}

\begin{exercise}\label{exercise:ch04-05}
一袋中含有 $a$ 只白球，$b$ 只黑球，从中摸出 $c$ 只（$c\leqslant a+b$），求摸出白球数 $\mu$ 的数学期望。
\end{exercise}

\begin{exercise}\label{exercise:ch04-06}
试求：(1) 为收集 $N$ 张赠券中的 $r$ 张所需购买的食品袋数 $\xi$ 的数学期望；(2) 为集齐水浒 $108$ 将，平均要购多少袋？（参照习题一第 $38$ 题）。
\end{exercise}

\begin{exercise}\label{exercise:ch04-07}
试证：若取非负整数值的随机变量 $\xi$ 的数学期望存在，则
\[
E\xi=\sum_{k=1}^{\infty}P\{\xi\geqslant k\}.
\]
\end{exercise}

\begin{exercise}\label{exercise:ch04-08}
若随机变量 $\xi$ 的分布函数为 $F(x)$，试证：
\[
E\xi=\int_0^{\infty}[1-F(x)]\,\mathrm{d}x-\int_{-\infty}^0 F(x)\,\mathrm{d}x.
\]
特别地，若 $\xi$ 取非负值，则
\[
E\xi=\int_0^{\infty}[1-F(x)]\,\mathrm{d}x.
\]
\end{exercise}

\begin{exercise}\label{exercise:ch04-09}
若随机变量 $\xi$ 服从 Laplace 分布，其密度函数为
\[
p(x)=\frac{1}{2\lambda}\mathrm{e}^{-|x-\mu|/\lambda},\quad -\infty<x<\infty,\ \lambda>0,
\]
试求 $E\xi$ 及 $D\xi$。
\end{exercise}

\begin{exercise}\label{exercise:ch04-10}
若分子的速度的分布密度函数由 Maxwell 分布律给出：
\[
p(s)=\sqrt{\frac{2}{\pi}}\cdot\frac{s^2}{\sigma^3}\exp\left(-\frac{s^2}{2\sigma^2}\right),\quad s>0,
\]
其中 $\sigma>0$ 是常数，试求分子的平均速度和平均动能（假定分子的质量等于 $m$）。
\end{exercise}

\begin{exercise}\label{exercise:ch04-11}
某城市共有 $N$ 辆汽车，车牌号从 $1$ 到 $N$，若随机地（可重复）记下 $n$ 辆车的车牌号，其最大号码为 $\xi$，求 $E\xi$。
\end{exercise}

\begin{exercise}\label{exercise:ch04-12}
若 $a\leqslant\xi\leqslant b$，试证：$D\xi\leqslant\dfrac{(b-a)^2}{4}$，并说明等式在何种情况下成立。
\end{exercise}

\begin{exercise}\label{exercise:ch04-13}
若 $\xi_1,\xi_2$ 相互独立，均服从 $N(\mu,\sigma^2)$，试证：
\[
E\max(\xi_1,\xi_2)=\mu+\frac{\sigma}{\sqrt{\pi}}.
\]
\end{exercise}

\begin{exercise}\label{exercise:ch04-14}
设 $f(x)\ (0\leqslant x<\infty)$ 是单调非降函数，且 $f(x)>0$。对随机变量 $\xi$，若 $Ef(|\xi|)<\infty$，则对任意 $x>0$，$P\{|\xi|\geqslant x\}\leqslant\dfrac{1}{f(x)}Ef(|\xi|)$。
\end{exercise}

\begin{exercise}\label{exercise:ch04-15}
若 $\xi_1,\xi_2,\cdots,\xi_n$ 为正的独立随机变量，服从相同分布，密度函数为 $p(x)$，试证
\[
E\left(\frac{\xi_1+\xi_2+\cdots+\xi_k}{\xi_1+\xi_2+\cdots+\xi_n}\right)=\frac{k}{n}.
\]
\end{exercise}

\begin{exercise}\label{exercise:ch04-16}
袋中装有 $N$ 只球，但其中白球数为随机变量，只知其数学期望为 $n$，试证从该袋中摸一球得到白球的概率为 $\dfrac{n}{N}$。
\end{exercise}

\begin{exercise}\label{exercise:ch04-17}
甲袋中装有 $a$ 只白球 $b$ 只黑球，乙袋中装有 $\alpha$ 只白球 $\beta$ 只黑球，现从甲袋中摸出 $c$ 只球（$c\leqslant a+b$）只球放入乙袋中，求从乙袋中再摸一球而为白球的概率。
\end{exercise}

\begin{exercise}\label{exercise:ch04-18}
袋中有 $a$ 只白球 $b$ 只黑球，每次摸出一球后总是放入一只白球，这样进行了 $n$ 次之后，再从袋中摸一只球，求它是白球的概率。
\end{exercise}

\begin{exercise}\label{exercise:ch04-19}
甲袋中有 $a$ 只白球 $b$ 只黑球，乙袋中有 $c$ 只白球 $d$ 只黑球，从两袋中各摸出一球，并交换放入另一袋中，这样做了 $n$ 次之后，再从甲袋中摸出一球，求这球是白球的概率。
\end{exercise}

\begin{exercise}\label{exercise:ch04-20}
现有 $n$ 个袋子，各装有 $a$ 只白球 $b$ 只黑球，先从第一个袋子中摸出一球，记下颜色后就把它放入第二个袋子中，再从第二个袋子中摸出一球，记下颜色后就把它放入第三个袋子中，照这样办法依次摸下去，最后从第 $n$ 个袋子中摸出一球并记下颜色，若在这 $n$ 次摸球中所摸得的白球的总数为 $S_n$，试求 $ES_n$。
\end{exercise}

\begin{exercise}\label{exercise:ch04-21}
在物理实验中，为测量某物体的重量，通常要重复测量多次，最后再把测量记录的平均值作为该物体的重量，试说明这样做的道理。
\end{exercise}

\begin{exercise}\label{exercise:ch04-22}
若 $\xi_1,\xi_2,\cdots,\xi_n$ 是独立随机变量，$D\xi_i=\sigma_i^2$，试找“权” $a_1,a_2,\cdots,a_n$（它们满足 $\sum_{i=1}^n a_i=1$），使 $\sum_{i=1}^n a_i\xi_i$ 的方差最小。
\end{exercise}

\begin{exercise}\label{exercise:ch04-23}
甲与乙依下列规则玩随机游戏：甲从装有 $i$ 号球 $i$ 个（$i=1,2,3,4,5$）的袋中随机摸出一球放入密盒中，让乙猜号。乙对甲的支付是他猜的号码与真正的号码之差的 (1) 平方；(2) 绝对值。试对这两种场合，讨论乙应采取的最佳策略。
\end{exercise}

\begin{exercise}\label{exercise:ch04-24}
某海港对停泊船只供给净水，初始价是每吨 $a$ 元，以后再供则要加 $50\%$ 的附加费；若用不完造成浪费则每吨加收资源费 $\dfrac{a}{4}$。设某轮船的净水用量是密度函数为 $p(x)$ 的随机变量，为节约其用水总开支试求其最佳首次供水量 $y$。
\end{exercise}

\begin{exercise}\label{exercise:ch04-25}
（Black--Scholes 期权定价公式）若股票价格 $S_T$ 服从对数正态分布，即 $\ln S_T\sim N\left(\ln s_t+\left(r-\dfrac{\sigma^2}{2}\right)(T-t),\ \sigma^2(T-t)\right)$，$t<T$。试证明该股票的敲定价为 $K$ 的买入期权的价格 $c_t=\mathrm{e}^{-r(T-t)}E[\max(S_T-K,0)]$ 满足如下 Black--Scholes 公式：$c_t=s_t\Phi(d_1)-K\mathrm{e}^{-r(T-t)}\Phi(d_2)$。

其中
\[
d_1=\frac{1}{\sigma\sqrt{T-t}}\left\{\ln\frac{s_t}{K}+\left(r+\frac{\sigma^2}{2}\right)(T-t)\right\},\qquad d_2=d_1-\sigma\sqrt{T-t}.
\]
\end{exercise}

\begin{exercise}\label{exercise:ch04-26}
Pareto 分布的密度函数为
\[
p(x)=\begin{cases}
rA^r\dfrac{1}{x^{r+1}},&x\geqslant A,\\
0,&x<A,
\end{cases}
\]
这里 $r>0,A>0$。试指出这分布具有 $p$ 阶矩，当且仅当 $p<r$。
\end{exercise}

\begin{exercise}\label{exercise:ch04-27}
若 $\xi$ 的密度函数为
\[
p(x)=\begin{cases}
\dfrac{1}{2|x|(\ln|x|)^2},&|x|>\mathrm{e},\\
0,&\text{其他},
\end{cases}
\]
试证对于任何 $\alpha>0$，$E|\xi|^\alpha=\infty$。
\end{exercise}

\begin{exercise}\label{exercise:ch04-28}
若 $\xi$ 服从 $N(\mu,\sigma^2)$，试求 $E|\xi-\mu|^k$，其中 $k$ 为正整数。
\end{exercise}

\begin{exercise}\label{exercise:ch04-29}
记 $a_k=E|\xi|^k$，若 $a_n<\infty$，试证 $\sqrt[k]{a_k}\leqslant\sqrt[k+1]{a_{k+1}},\ k=1,2,\cdots,n-1$。
\end{exercise}

\begin{exercise}\label{exercise:ch04-30}
设随机变量 $\xi_1,\xi_2,\cdots,\xi_{m+n}\ (n>m)$ 是独立的，有相同的分布并且有有限的方差，试求 $S=\xi_1+\cdots+\xi_n$ 与 $T=\xi_{m+1}+\xi_{m+2}+\cdots+\xi_{m+n}$ 两和之间的相关系数。
\end{exercise}

\begin{exercise}\label{exercise:ch04-31}
若 $(\xi,\eta)$ 的密度函数为
\[
p(x,y)=\begin{cases}
\dfrac{1}{\pi},&x^2+y^2\leqslant 1,\\
0,&x^2+y^2>1,
\end{cases}
\]
试验证：$\xi$ 与 $\eta$ 不相关，但它们不独立。
\end{exercise}

\begin{exercise}\label{exercise:ch04-32}
某人写好 $n$ 封信，又写好 $n$ 只信封，在黑暗中把每封信随意放入某一信封中，试求放对的信封数 $\mu$ 的数学期望及方差。
\end{exercise}

\begin{exercise}\label{exercise:ch04-33}
设随机变量 $\zeta\sim N(0,1)$，记 $A=\{|\zeta|>1\}$，$B=\{|\zeta|>2\}$，试求 $\mathbf{1}_A$ 及 $\mathbf{1}_B$ 的概率分布列联表，数学期望，方差，它们的相关系数以及 $D(\mathbf{1}_A+\mathbf{1}_B)$。
\end{exercise}

\begin{exercise}\label{exercise:ch04-34}
若 $\xi,\eta$ 服从二元正态分布，$E\xi=a,D\xi=1,E\eta=b,D\eta=1$，证明：$\xi$ 与 $\eta$ 的相关系数 $r=\cos q\pi$，其中 $q=P\{(\xi-a)(\eta-b)<0\}$。
\end{exercise}

\begin{exercise}\label{exercise:ch04-35}
设 $(\xi,\eta)$ 服从二元正态分布，$E\xi=E\eta=0,D\xi=D\eta=1,r_{\xi\eta}=\rho$，试证
\[
E\max(\xi,\eta)=\sqrt{\frac{1-\rho}{\pi}}.
\]
\end{exercise}

\begin{exercise}\label{exercise:ch04-36}
甲袋中装有 $5$ 只白球，$7$ 只黑球，$3$ 只红球，乙袋中装有 $4$ 只白球，$4$ 只黑球，$7$ 只红球，试问从哪一个袋中取出一只球有较大不肯定性？
\end{exercise}

\begin{exercise}\label{exercise:ch04-37}
试求几何分布的熵。
\end{exercise}

\begin{exercise}\label{exercise:ch04-38}
试求二项分布的熵。
\end{exercise}

\begin{exercise}\label{exercise:ch04-39}
若以 $\alpha$ 及 $\beta$ 分别记二进位信道的输入及输出，已知 $P\{\alpha=1\}=p$，$P\{\alpha=0\}=1-p$，$P\{\beta=1\mid\alpha=1\}=q$，$P\{\beta=1\mid\alpha=0\}=1-q$，$P\{\beta=1\mid\alpha=0\}=r$，$P\{\beta=0\mid\alpha=0\}=1-r$，试求输出中含有输入的信息量。
\end{exercise}

\begin{exercise}\label{exercise:ch04-40}
在 $12$ 只金属球中，混有一只假球，并且不知道它是比真球轻还是重，用没有砝码的天平来称这些球。(1) 试问至少需要称多少次才能查出这个假球并确定它是比真球轻或重。(2) 给出一种称球方案。
\end{exercise}

\begin{exercise}\label{exercise:ch04-41}
试用母函数法求 Pascal 分布的数学期望及方差。
\end{exercise}

\begin{exercise}\label{exercise:ch04-42}
设 $\xi$ 是一个母函数为 $P(s)$ 的随机变量，试求下列各概率所对应的母函数：

(1) $P\{\xi>n\}$；\qquad (2) $P\{\xi=2n\}$。
\end{exercise}

\begin{exercise}\label{exercise:ch04-43}
在 Bernoulli 试验中，若试验次数 $\nu$ 是随机变量，试证成功的次数与失败的次数这两个随机变量独立的充要条件是 $\nu$ 服从 Poisson 分布。
\end{exercise}

\begin{exercise}\label{exercise:ch04-44}
设 $\{\xi_k\}$ 是一串独立的整值随机变量序列，具有相同概率分布，考虑和 $\eta=\xi_1+\xi_2+\cdots+\xi_\nu$，其中 $\nu$ 是随机变量，它与 $\{\xi_k\}$ 相互独立，试用 (1) 母函数法；(2) 直接计算证明
\[
E\eta=E\nu\cdot E\xi_k,\qquad D\eta=E\nu\cdot D\xi_k+D\nu\cdot(E\xi_k)^2.
\]
\end{exercise}

\begin{exercise}\label{exercise:ch04-45}
某公共汽车站在 $[0,t]$ 中来到的乘客批数 $\mu$ 服从参数为 $\lambda t$ 的 Poisson 分布，而每批来到的乘客数是随机变量，来 $n$ 个的概率为 $p_n,\ n=0,1,2,\cdots$。试求 $[0,t]$ 中来到乘客数 $\eta$ 的母函数及数学期望。
\end{exercise}

\begin{exercise}\label{exercise:ch04-46}
试用母函数法证明二项分布、Poisson 分布与 Pascal 分布的再生性。
\end{exercise}

\begin{exercise}\label{exercise:ch04-47}
若分布函数 $F(x)=1-F(-x+0)$ 成立，则称它是对称的；试证分布函数对称的充要条件是它的特征函数是实的偶函数。
\end{exercise}

\begin{exercise}\label{exercise:ch04-48}
试求 $[0,1]$ 均匀分布的特征函数。
\end{exercise}

\begin{exercise}\label{exercise:ch04-49}
一般 Cauchy 分布的密度函数为 $p(x)=\dfrac{1}{\pi}\cdot\dfrac{\lambda}{\lambda^2+(x-\mu)^2},\ \lambda>0$，试证它的特征函数为 $\mathrm{e}^{\mathrm{i}\mu t-\lambda|t|}$，利用这个结果证明 Cauchy 分布的再生性。
\end{exercise}

\begin{exercise}\label{exercise:ch04-50}
若随机变量 $\xi$ 服从 Cauchy 分布，$\mu=0,\lambda=1$，而 $\eta=\xi$，试证关于特征函数成立着
\[
f_{\xi+\eta}(t)=f_{\xi}(t)\cdot f_{\eta}(t),
\]
但是 $\xi$ 与 $\eta$ 并不独立。
\end{exercise}

\begin{exercise}\label{exercise:ch04-51}
设 $\xi_1,\xi_2,\cdots,\xi_n$ 相互独立且均服从同一 Cauchy 分布，试证：$\dfrac{1}{n}(\xi_1+\xi_2+\cdots+\xi_n)$ 与 $\xi_1$ 同分布。
\end{exercise}

\begin{exercise}\label{exercise:ch04-52}
若 $\xi\sim N(\mu,\sigma^2)$，试用特征函数法求 $E(\xi-\mu)^k$。
\end{exercise}

\begin{exercise}\label{exercise:ch04-53}
求证：对于任何实值特征函数 $f(t)$，以下两个不等式成立：
\[
1-f(2t)\leqslant 4(1-f(t)),
\]
\[
1+f(2t)\geqslant 2(f(t))^2.
\]
\end{exercise}

\begin{exercise}\label{exercise:ch04-54}
求证，如果 $f(t)$ 是相应于分布函数 $F(x)$ 的特征函数，则对于任何 $x$ 值恒成立
\[
\lim_{T\to\infty}\frac{1}{2T}\int_{-T}^T f(t)\mathrm{e}^{-\mathrm{i}tx}\,\mathrm{d}t=F(x+0)-F(x-0).
\]
\end{exercise}

\begin{exercise}\label{exercise:ch04-55}
随机变量 $\xi$ 的特征函数为 $f(t)$，且它的 $n$ 阶矩存在，令
\[
\chi_k=\frac{1}{\mathrm{i}^k}\left[\frac{\mathrm{d}^k}{\mathrm{d}t^k}\log f(t)\right]_{t=0},\qquad k\leqslant n,
\]
称 $\chi_k$ 为随机变量 $\xi$ 的 $k$ 阶\textbf{半不变量}。

(1) 试证 $\eta=\xi+b$（$b$ 是常数）的 $k\,(k>1)$ 阶半不变量等于 $\chi_k$。

(2) 试求出半不变量与原点矩之间的关系式。
\end{exercise}

\begin{exercise}\label{exercise:ch04-56}
若随机向量 $(\xi,\eta)$ 服从二元正态分布 $N(\mu_1,\mu_2,\sigma_1^2,\sigma_2^2,\rho)$，试写出：

(1) $(\xi,\eta)$ 的特征函数；

(2) $a\xi+b\eta$ 的密度函数；

(3) 当 $\xi=x$ 时 $\eta$ 的条件密度函数，并讨论 $\rho=-1,\rho=1$ 及 $\rho=0$ 等特殊情况下结果的概率含意。
\end{exercise}

\begin{exercise}\label{exercise:ch04-57}
若 $\xi_1,\xi_2,\cdots,\xi_n$ 相互独立，均服从 $N(0,1)$，而
\[
\eta_1=\sum_{k=1}^n a_k\xi_k,\qquad \eta_2=\sum_{k=1}^n b_k\xi_k,
\]
试证 $\eta_1$ 与 $\eta_2$ 独立的充要条件为 $\sum_{k=1}^n a_k b_k=0$。
\end{exercise}

\begin{exercise}\label{exercise:ch04-58}
设 $\xi_1,\xi_2,\cdots,\xi_n$ 相互独立，具有相同分布 $N(\mu,\sigma^2)$，试求 $\boldsymbol{\xi}=\begin{pmatrix}\xi_1\\\xi_2\\\vdots\\\xi_n\end{pmatrix}$ 的分布，并写出它的数学期望及协方差矩阵。再求 $\bar{\xi}=\dfrac{1}{n}\sum_{i=1}^n\xi_i$ 的分布密度。
\end{exercise}

\begin{exercise}\label{exercise:ch04-59}
若 $(\xi,\eta)$ 服从 $N(\mu_1,\mu_2,\sigma_1^2,\sigma_2^2,\rho)$，而
\[
U=a\xi+b\eta,\qquad V=c\xi+d\eta.
\]
(1) 试求 $U$ 与 $V$ 的数学期望，方差及相关系数；

(2) 写出 $(U,V)$ 的分布；

(3) 讨论：何种情况下，$(U,V)$ 退化为一维分布；何种情况下，$U$ 与 $V$ 独立。
\end{exercise}

\begin{exercise}\label{exercise:ch04-60}
（Fisher 引理）若 $X_1,X_2,\cdots,X_n$ 相互独立，均服从 $N(\mu,\sigma^2)$，记
\[
\bar{X}=\frac{1}{n}\sum_{i=1}^n X_i,\qquad S_n^2=\frac{1}{n}\sum_{i=1}^n(X_i-\bar{X})^2.
\]
试证：

(1) $\bar{X}$ 与 $S_n^2$ 相互独立；

(2) $\bar{X}\sim N\left(\mu,\dfrac{\sigma^2}{n}\right)$；

(3) $\dfrac{nS_n^2}{\sigma^2}\sim\chi_{n-1}^2$。
\end{exercise}

\end{document}